\documentclass[11pt, a4paper]{article}
\usepackage[margin=1in]{geometry}
\usepackage{amsmath}
\usepackage{graphicx}
\usepackage{booktabs}
\usepackage{hyperref}
\usepackage{xcolor}
\usepackage{listings}
\usepackage{tcolorbox}
\usepackage{enumitem}

\tcbuselibrary{skins, breakable, listings}

% Code listing style
\lstdefinestyle{bash}{
    language=bash,
    basicstyle=\ttfamily\small,
    backgroundcolor=\color{gray!10},
    frame=single,
    framerule=0pt,
    breaklines=true,
    columns=fullflexible,
    keepspaces=true,
    showstringspaces=false,
    commentstyle=\color{green!50!black},
    keywordstyle=\color{blue!70!black},
    stringstyle=\color{red!70!black},
    morekeywords={pip, python, kaggle, mkdir, cd, wget, git, unzip},
}

\lstdefinestyle{python}{
    language=Python,
    basicstyle=\ttfamily\small,
    backgroundcolor=\color{blue!3},
    frame=single,
    framerule=0pt,
    breaklines=true,
    columns=fullflexible,
    keepspaces=true,
    showstringspaces=false,
    commentstyle=\color{green!50!black},
    keywordstyle=\color{blue!70!black}\bfseries,
    stringstyle=\color{red!70!black},
}

\newtcolorbox{important}[1][]{
    colback=red!5!white, colframe=red!75!black,
    fonttitle=\bfseries, title={#1}, breakable
}
\newtcolorbox{tip}[1][]{
    colback=green!5!white, colframe=green!50!black,
    fonttitle=\bfseries, title={#1}, breakable
}
\newtcolorbox{step}[1][]{
    colback=blue!5!white, colframe=blue!50!black,
    fonttitle=\bfseries, title={#1}, breakable
}

\title{\textbf{Complete Kaggle Training Guide}\\[0.3em]
\Large DP-Guided Decoupled Mask-and-Fill Pipeline\\
Step-by-Step from Zero to Trained Models}
\author{Team: Mr.A}
\date{\today}

\begin{document}
\maketitle
\tableofcontents
\newpage

% ============================================================
\section{Prerequisites}
% ============================================================

\subsection{What You Need}
\begin{itemize}
    \item A free Kaggle account: \url{https://www.kaggle.com}
    \item Phone-verified account (required for GPU access)
    \item Kaggle provides \textbf{30 hours/week} of free GPU (T4 or P100)
    \item Internet access enabled in the notebook settings
\end{itemize}

\subsection{Hardware Available on Kaggle}
\begin{table}[h]
\centering
\begin{tabular}{lll}
\toprule
\textbf{Accelerator} & \textbf{VRAM} & \textbf{Best For} \\
\midrule
GPU T4 $\times$ 2 & 16 GB each & QLoRA training (our pipeline) \\
GPU P100 $\times$ 1 & 16 GB & Inference + evaluation \\
TPU v3-8 & 128 GB HBM & Full fine-tuning (not needed for us) \\
\bottomrule
\end{tabular}
\end{table}

\begin{important}[GPU Quota]
Kaggle gives 30 GPU hours per week. Our full pipeline takes approximately \textbf{8--12 hours} on a T4 GPU, well within the weekly limit. If you run out, wait for the weekly reset (Saturday UTC) or use a second Kaggle account.
\end{important}

% ============================================================
\section{Step 1: Create a New Kaggle Notebook}
% ============================================================

\begin{step}[Instructions]
\begin{enumerate}
    \item Go to \url{https://www.kaggle.com/code}
    \item Click \textbf{``+ New Notebook''}
    \item In the right panel, configure:
    \begin{itemize}
        \item \textbf{Language}: Python
        \item \textbf{Accelerator}: GPU T4 $\times$ 2
        \item \textbf{Internet}: \textbf{ON} (toggle this --- required for downloading models)
        \item \textbf{Persistence}: Files (so outputs survive session restarts)
    \end{itemize}
    \item Rename the notebook to: \texttt{privacy-anonymization-v8}
\end{enumerate}
\end{step}

\begin{important}[Internet Must Be ON]
If Internet is OFF, you cannot download models from Hugging Face or install packages. Toggle it ON in \textbf{Settings $\rightarrow$ Internet $\rightarrow$ On}. You may need to verify your phone number first.
\end{important}

% ============================================================
\section{Step 2: Install Dependencies}
% ============================================================

In the \textbf{first cell} of your notebook, paste:

\begin{lstlisting}[style=bash]
# Cell 1: Install all dependencies
!pip install -q -U \
    "transformers>=4.40.0" \
    datasets \
    evaluate \
    accelerate \
    "peft>=0.10.0" \
    bitsandbytes \
    rouge_score \
    sacrebleu \
    sentencepiece \
    scipy \
    scikit-learn \
    bert_score \
    Faker \
    opacus \
    presidio-analyzer \
    presidio-anonymizer \
    gradio

# Verify GPU
import torch
print(f"GPU available: {torch.cuda.is_available()}")
print(f"GPU name: {torch.cuda.get_device_name(0)}")
print(f"VRAM: {torch.cuda.get_device_properties(0).total_mem / 1e9:.1f} GB")
\end{lstlisting}

\textbf{Expected output:}
\begin{lstlisting}[style=bash]
GPU available: True
GPU name: Tesla T4
VRAM: 15.8 GB
\end{lstlisting}

% ============================================================
\section{Step 3: Upload the Script}
% ============================================================

You have two options:

\subsection{Option A: Copy-Paste (Recommended)}

\begin{enumerate}
    \item Open \texttt{script.py} from your local project folder.
    \item Copy the entire contents.
    \item In Kaggle, create a new \textbf{Code cell} and paste everything.
    \item The script is designed with \texttt{\# \%\%} cell markers --- Kaggle will recognize these as cell boundaries if you use ``Convert to cells'' option.
\end{enumerate}

\subsection{Option B: Upload as Dataset}

\begin{enumerate}
    \item Go to \url{https://www.kaggle.com/datasets} $\rightarrow$ \textbf{+ New Dataset}
    \item Upload \texttt{script.py} as a file
    \item In your notebook, add the dataset: \textbf{+ Add Data $\rightarrow$ Your Datasets}
    \item Then in a cell:
\begin{lstlisting}[style=python]
# Load the script from the dataset
exec(open("/kaggle/input/your-dataset-name/script.py").read())
\end{lstlisting}
\end{enumerate}

\begin{tip}[Recommendation]
Option A is simpler and lets you edit code directly in Kaggle. The script has \texttt{\# \%\%} markers that work as cell separators in Kaggle notebooks, giving you a clean multi-cell layout.
\end{tip}

% ============================================================
\section{Step 4: Configure for Kaggle Environment}
% ============================================================

Before running, modify the \texttt{Config} class for Kaggle paths:

\begin{lstlisting}[style=python]
# Cell 2: Override Config for Kaggle
class Config:
    MODE = "ALL"

    # Model backbones
    CENSOR_CHECKPOINT = "microsoft/deberta-v3-base"
    HALLUC_CHECKPOINT = "google/flan-t5-large"
    ZEROSHOT_CHECKPOINT = "google/flan-t5-large"  # Use same as halluc to save VRAM

    # Kaggle paths (use /kaggle/working for outputs)
    OUTPUT_ROOT = "/kaggle/working/privacy_project_v8"
    MODEL_CENSOR_DIR = "/kaggle/working/privacy_project_v8/censor"
    MODEL_HALLUC_DIR = "/kaggle/working/privacy_project_v8/hallucinator"
    EVAL_DIR = "/kaggle/working/privacy_project_v8/evaluation"

    # Hyperparameters (tuned for T4 16GB VRAM)
    MAX_LEN = 256
    NER_MAX_LEN = 256
    BATCH_SIZE = 4            # Smaller batch for T4
    GRAD_ACCUMULATION = 8     # Effective batch = 32
    EPOCHS = 3
    LR_CENSOR = 3e-5
    LR_HALLUC = 1e-4
    WARMUP_RATIO = 0.06

    # QLoRA
    LORA_R = 32
    LORA_ALPHA = 64
    LORA_DROPOUT = 0.05

    # DP-SGD
    ENABLE_DP = True
    DP_EPSILON = 8.0
    DP_DELTA = None
    DP_MAX_GRAD_NORM = 1.0

    # Generation
    GEN_MAX_TOKENS = 200
    GEN_NUM_BEAMS = 4
    GEN_TEMPERATURE = 0.7
    GEN_TOP_K = 50
    GEN_TOP_P = 0.9
    GEN_REPETITION_PENALTY = 1.2

    # Data -- USE SUBSAMPLING FOR FIRST RUN
    TARGET_SAMPLE_COUNT = 10000   # Start with 10K to verify everything works
    TEST_SIZE = 0.05
    NUM_EVAL_SAMPLES = 100

    ENTITY_TYPES = [
        "PERSON", "LOC", "ORG", "DATE", "PHONE", "EMAIL",
        "SSN", "CREDIT_CARD", "ADDRESS", "IP_ADDRESS",
    ]
\end{lstlisting}

\begin{important}[VRAM Management on T4]
The T4 has 16GB VRAM. Key settings to avoid OOM (Out of Memory):
\begin{itemize}[nosep]
    \item \texttt{BATCH\_SIZE = 4} (not 8 or 16)
    \item \texttt{GRAD\_ACCUMULATION = 8} (compensates for small batch)
    \item Always use \texttt{load\_in\_4bit=True} (QLoRA)
    \item Use \texttt{torch.bfloat16} compute dtype
    \item If OOM still occurs, reduce \texttt{MAX\_LEN} to 192 or \texttt{LORA\_R} to 16
\end{itemize}
\end{important}

% ============================================================
\section{Step 5: Run the Data Pipeline}
% ============================================================

\begin{lstlisting}[style=python]
# Cell 3: Load and prepare data
# This is already in the script - just run the data pipeline section

# After running, you should see output like:
# Loading AI4Privacy dataset...
# Raw dataset size: 500000
# After PII filter: 487231
# English subset: 148923
# Subsampled to 10000
# Partitions -- Baseline: 950 | Censor (A): 4512 | Hallucinator (B): 4513 | Test: 500
\end{lstlisting}

\begin{tip}[Quick Test First]
Set \texttt{TARGET\_SAMPLE\_COUNT = 10000} for your first run. This completes in $\sim$2 hours and lets you verify everything works. Once confirmed, set it to \texttt{None} for the full 150K English dataset ($\sim$8--10 hours).
\end{tip}

% ============================================================
\section{Step 6: Train the Censor (Model I)}
% ============================================================

\begin{lstlisting}[style=python]
# Cell 4: Train Censor
# This runs automatically from the script's training section
# Expected output:

# ============================================================
# TRAINING CENSOR (Flan-T5-Large + QLoRA, Seq2Seq)
# ============================================================
# trainable params: 19,922,944 || all params: 803,430,400 || trainable%: 2.48
# {'loss': 2.1534, 'step': 25}
# {'loss': 1.4823, 'step': 50}
# {'loss': 0.9127, 'step': 75}
# ...
# Loss curve saved: /kaggle/working/.../censor/loss_curve.png
# Censor training complete.
\end{lstlisting}

\textbf{Expected training time:} $\sim$45 minutes (10K samples) to $\sim$4 hours (full dataset) on T4.

% ============================================================
\section{Step 7: Train the Hallucinator (Model II)}
% ============================================================

\begin{lstlisting}[style=python]
# Cell 5: Train Hallucinator
# Runs automatically after Censor training
# Expected output:

# ============================================================
# TRAINING HALLUCINATOR (Flan-T5-Large + QLoRA)
# ============================================================
# trainable params: 19,922,944 || all params: 803,430,400 || trainable%: 2.48
# {'loss': 3.2145, 'step': 25}
# {'loss': 1.8934, 'step': 50}
# ...
# Loss curve saved: /kaggle/working/.../hallucinator/loss_curve.png
# Hallucinator training complete.
\end{lstlisting}

\textbf{Expected training time:} Same as Censor ($\sim$45 minutes for 10K, $\sim$4 hours for full).

% ============================================================
\section{Step 8: Run Evaluation}
% ============================================================

The evaluation section runs automatically after training:

\begin{lstlisting}[style=python]
# Cell 6: Evaluation runs automatically
# Expected output:

# ============================================================
# COMPREHENSIVE EVALUATION
# ============================================================
#
# --- Approach 2: Presidio Baseline ---
#   Leakage: {'exact_leakage_rate': 0.03, ...}
#   Utility: {'BLEU': 28.4, 'ROUGE-L': 45.2, ...}
#
# --- Approach 3: Zero-Shot LLM ---
#   Leakage: {'exact_leakage_rate': 0.08, ...}
#   Utility: {'BLEU': 32.1, 'ROUGE-L': 51.7, ...}
#
# --- Approach 4: Decoupled Mask-and-Fill ---
#   Leakage: {'exact_leakage_rate': 0.005, ...}
#   Utility: {'BLEU': 38.2, 'ROUGE-L': 58.3, ...}
#
# COMPARISON SUMMARY
# | Method    | Exact_Leakage% | BLEU  | BERTScore_F1 |
# |-----------|----------------|-------|--------------|
# | Presidio  |           3.00 | 28.40 |        82.30 |
# | ZeroShot  |           8.00 | 32.10 |        85.10 |
# | Decoupled |           0.50 | 38.20 |        91.40 |
\end{lstlisting}

% ============================================================
\section{Step 9: Save and Download Results}
% ============================================================

\begin{lstlisting}[style=python]
# Cell 7: Package all outputs for download
import shutil, os

output_dir = "/kaggle/working/privacy_project_v8"

# Zip everything
shutil.make_archive(
    "/kaggle/working/final_results",
    "zip",
    output_dir
)

print("Zip created at: /kaggle/working/final_results.zip")
print(f"Size: {os.path.getsize('/kaggle/working/final_results.zip') / 1e6:.1f} MB")

# List all output files
for root, dirs, files in os.walk(output_dir):
    for f in files:
        path = os.path.join(root, f)
        size = os.path.getsize(path) / 1e6
        print(f"  {path} ({size:.1f} MB)")
\end{lstlisting}

\begin{step}[Download Your Results]
\begin{enumerate}
    \item In Kaggle, click \textbf{``Save Version''} (top right) $\rightarrow$ \textbf{``Save \& Run All''}
    \item Wait for execution to complete (check the ``Versions'' tab)
    \item Go to the Output tab of your notebook
    \item Download \texttt{final\_results.zip}
    \item This contains: trained model adapters, loss curves, evaluation CSVs, comparison tables
\end{enumerate}
\end{step}

% ============================================================
\section{Step 10: Full Production Run}
% ============================================================

Once the quick test (Step 5--9) succeeds, run the full pipeline:

\begin{lstlisting}[style=python]
# Change these in Config:
Config.TARGET_SAMPLE_COUNT = None   # Use FULL dataset
Config.NUM_EVAL_SAMPLES = 200      # More thorough evaluation
Config.EPOCHS = 3                  # Keep at 3
\end{lstlisting}

\begin{table}[h]
\centering
\begin{tabular}{lrr}
\toprule
\textbf{Stage} & \textbf{Quick Test (10K)} & \textbf{Full Run (150K)} \\
\midrule
Data loading & 2 min & 5 min \\
Censor training & 45 min & 4 hrs \\
Hallucinator training & 45 min & 4 hrs \\
Presidio baseline & 5 min & 10 min \\
Zero-shot baseline & 15 min & 30 min \\
Evaluation & 10 min & 30 min \\
\midrule
\textbf{Total} & \textbf{$\sim$2 hrs} & \textbf{$\sim$9 hrs} \\
\bottomrule
\end{tabular}
\caption{Expected execution times on Kaggle T4 GPU.}
\end{table}

% ============================================================
\section{Troubleshooting}
% ============================================================

\subsection{Common Issues}

\begin{table}[h]
\centering
\small
\begin{tabular}{p{4cm}p{4cm}p{4cm}}
\toprule
\textbf{Error} & \textbf{Cause} & \textbf{Fix} \\
\midrule
\texttt{CUDA out of memory} & Batch too large for T4 & Set \texttt{BATCH\_SIZE=2}, \texttt{GRAD\_ACCUMULATION=16} \\
\midrule
\texttt{RuntimeError: Expected all tensors on same device} & Mixed CPU/GPU tensors & Add \texttt{.to(device)} to all inputs before generation \\
\midrule
\texttt{Connection error} downloading model & Internet not enabled & Settings $\rightarrow$ Internet $\rightarrow$ ON \\
\midrule
\texttt{No module named 'bitsandbytes'} & Deps not installed & Re-run Cell 1 (install cell) \\
\midrule
Training loss stuck at high value & LR too high or data issue & Lower LR to \texttt{5e-5}; verify data pipeline output \\
\midrule
\texttt{Disk quota exceeded} & Kaggle has 20GB working dir & Delete intermediate checkpoints: \\
& & \texttt{!rm -rf /kaggle/working/*/checkpoint-*} \\
\midrule
Session timeout (12hr limit) & Kaggle kills sessions at 12hrs & Use \texttt{Save \& Run All} for unattended execution \\
\bottomrule
\end{tabular}
\end{table}

\subsection{VRAM Optimization Checklist}

If you encounter OOM errors, apply these in order:

\begin{enumerate}
    \item \textbf{Reduce batch size}: \texttt{BATCH\_SIZE = 2} (increase \texttt{GRAD\_ACCUMULATION} proportionally)
    \item \textbf{Reduce sequence length}: \texttt{MAX\_LEN = 192} or \texttt{128}
    \item \textbf{Reduce LoRA rank}: \texttt{LORA\_R = 16} (from 32)
    \item \textbf{Use gradient checkpointing}: Add \texttt{gradient\_checkpointing=True} to training args
    \item \textbf{Clear cache between models}:
\begin{lstlisting}[style=python]
import gc, torch
gc.collect()
torch.cuda.empty_cache()
\end{lstlisting}
    \item \textbf{Downgrade model}: Use \texttt{google/flan-t5-base} (250M) instead of \texttt{flan-t5-large} (780M)
\end{enumerate}

% ============================================================
\section{Notebook Cell Structure (Copy-Paste Reference)}
% ============================================================

For the cleanest Kaggle experience, organize your notebook as these cells:

\begin{table}[h]
\centering
\small
\begin{tabular}{clll}
\toprule
\textbf{Cell} & \textbf{Type} & \textbf{Content} & \textbf{Run Time} \\
\midrule
1 & Code & Install dependencies & 2 min \\
2 & Code & Imports + Config (override for Kaggle) & 10 sec \\
3 & Code & Data pipeline (\texttt{load\_ai4privacy} + \texttt{prepare\_data}) & 3 min \\
4 & Code & Train Censor (\texttt{train\_censor\_seq2seq}) & 45 min \\
5 & Code & Train Hallucinator (\texttt{train\_hallucinator}) & 45 min \\
6 & Code & Presidio baseline (\texttt{run\_presidio\_baseline}) & 5 min \\
7 & Code & Zero-shot baseline (\texttt{run\_zeroshot\_baseline}) & 15 min \\
8 & Code & Decoupled inference (\texttt{run\_decoupled\_pipeline}) & 10 min \\
9 & Code & Evaluation metrics + comparison table & 5 min \\
10 & Code & Save \& zip results & 1 min \\
\bottomrule
\end{tabular}
\caption{Recommended cell layout for Kaggle notebook.}
\end{table}

% ============================================================
\section{Verifying Your Results}
% ============================================================

After the full run, you should have these files in \texttt{/kaggle/working/privacy\_project\_v8/}:

\begin{lstlisting}[style=bash]
privacy_project_v8/
  censor/
    best/                    # Trained Censor adapter weights
      adapter_config.json
      adapter_model.safetensors
    loss_curve.png           # Training loss plot
  hallucinator/
    best/                    # Trained Hallucinator adapter weights
      adapter_config.json
      adapter_model.safetensors
    loss_curve.png
  evaluation/
    decoupled_results.csv    # Per-example anonymization results
    comparison_summary.csv   # Method comparison table
    all_results.json         # All metrics in JSON
\end{lstlisting}

\subsection{Sanity Checks}
\begin{enumerate}
    \item \textbf{Loss curves}: Should show decreasing trend; final loss $< 1.0$ for both models.
    \item \textbf{Entity leakage rate}: Decoupled method should have $< 1\%$ exact leakage.
    \item \textbf{BERTScore}: Should be $> 0.85$ for the decoupled method.
    \item \textbf{Sample outputs}: Open \texttt{decoupled\_results.csv} and manually verify that:
    \begin{itemize}
        \item Names are replaced with different realistic names
        \item Locations are replaced with different real locations
        \item Grammar is preserved
        \item Non-entity text is unchanged
    \end{itemize}
\end{enumerate}

\begin{tip}[Save Your Notebook Version]
Always click \textbf{``Save Version'' $\rightarrow$ ``Save \& Run All (Commit)''} before closing. This creates a permanent snapshot with all outputs. You can download outputs from the ``Output'' tab of any committed version, even weeks later.
\end{tip}

\end{document}
